\chapter{Graves' Eye Disease}
\index{Graves' Eye Disease}
\label{sec:Graves' Eye Disease}

\section{Overview}
\begin{itemize}[noitemsep]
  \item Graves' eye disease, also called thyroid eye disease or Graves' orbitopathy, is an autoimmune inflammation of the tissues behind the eye
  \item It affects about 25--50\% of people with Graves' disease (overactive thyroid) and can occur with any thyroid condition
  \item Symptoms range from mild dryness and eye prominence to double vision and, rarely, sight-threatening optic nerve compression
\end{itemize}
\section{Causes}
This condition happens because of an autoimmune reaction in which antibodies stimulate fibroblasts and immune cells in the eye socket (orbit), causing inflammation, swelling, and fat and muscle enlargement. The expanding tissue pushes the eye forward (proptosis). Smoking, thyroid dysfunction, and radiation therapy for thyroid disease can worsen the condition.
\section{Risk Factors}
\begin{itemize}[noitemsep]
  \item Graves' disease (hyperthyroidism)
  \item Female sex and middle age
  \item Cigarette smoking (strongly increases severity)
  \item Radioactive iodine treatment that causes fluctuating thyroid hormone levels
  \item Family history of autoimmune thyroid disease
  \item High thyroid antibody levels
\end{itemize}
\section{Signs and Symptoms}
\subsection{Common symptoms}
\begin{itemize}[noitemsep]
  \item Bulging eyes (proptosis)
  \item Dry, gritty, or watery eyes
  \item Redness, puffiness, and swelling around the eyes
  \item Sensitivity to light and double vision (diplopia)
  \item Difficulty fully closing the eyelids
  \item Pain behind the eyes worsened by eye movement
\end{itemize}
\subsection{Emergency warning signs}
\begin{itemize}[noitemsep]
  \item Sudden vision loss, severe double vision, or intense eye pain
  \item Inability to close the eyes causing corneal exposure
  \item Signs of optic nerve compression, such as reduced color vision or vision loss --- these require urgent specialist care
\end{itemize}
\section{Progression and Stages}
The disease follows an active (inflammatory) phase lasting 6 months to 2 years, followed by a stable (inactive) phase. Progression ranges from mild to severe. The majority of cases are mild, but a minority progress to moderate or severe disease with significant proptosis, diplopia, or optic neuropathy. Treatment in the active phase aims to reduce inflammation; rehabilitation is deferred to the stable phase.
\section{Complications}
\begin{itemize}[noitemsep]
  \item Optic nerve compression leading to vision loss
  \item Corneal exposure, dryness, and ulceration
  \item Permanent double vision
  \item Disfiguring proptosis
  \item Facial and psychological distress from appearance changes
\end{itemize}
\section{Diagnosis}
\begin{itemize}[noitemsep]
  \item Clinical examination, including measurement of eye protrusion (exophthalmometry) and eye movement
  \item Thyroid function tests (TSH, free T4, T3) and thyroid antibodies (TSI, TRAb)
  \item Orbital CT or MRI to assess tissue swelling and muscle involvement
  \item Visual acuity, color vision, and visual field testing to check for optic nerve compression
  \item Optical coherence tomography (OCT) of the optic nerve in severe cases
\end{itemize}
\section{Prevention}
\begin{itemize}[noitemsep]
  \item Quit smoking, as this is the strongest modifiable risk factor
  \item Stable control of thyroid hormone levels
  \item Avoidance of hypothyroidism after thyroid treatment
  \item Early referral to an eye specialist when eye symptoms develop
\end{itemize}
\section{Treatment Options}
\begin{itemize}[noitemsep]
  \item Selenium supplements and lubricating eye drops for mild disease
  \item Corticosteroids (oral or intravenous) for active moderate-to-severe inflammation
  \item Immunomodulators (e.g., rituximab, tocilizumab) or teprotumumab for active disease
  \item Orbital radiation therapy in selected cases
  \item Surgery for stable disease: orbital decompression for proptosis, strabismus surgery for diplopia, eyelid surgery
  \item Management of thyroid status by an endocrinologist
\end{itemize}
\section{Living with It}
\begin{itemize}[noitemsep]
  \item Protect eyes with sunglasses and lubricants, especially outdoors
  \item Elevate the head while sleeping to reduce periorbital swelling
  \item Seek support for appearance concerns and connect with support groups
  \item Attend coordinated endocrinology and ophthalmology follow-up
\end{itemize}
\section{Outlook}
Most cases of Graves' eye disease are mild and self-limited, and the active phase eventually burns out. With modern therapy, severe disease is usually controlled, and surgery restores function and appearance in most. Vision loss from optic neuropathy is permanent if untreated but is rare when managed promptly.
